% Generated draft: requires mathematical and source review.
\documentclass[11pt]{article}
\usepackage[margin=25mm]{geometry}
\usepackage{fontspec}
\setmainfont{Latin Modern Roman}
\usepackage{amsmath,amssymb,mathtools,unicode-math}
\setmathfont{Latin Modern Math}
\usepackage{enumitem,needspace,xurl,hyperref}
\setlength{\parindent}{0pt}
\setlength{\parskip}{5pt}
\setlength{\emergencystretch}{3em}
\newcommand{\N}{\mathbb{N}}
\newcommand{\Z}{\mathbb{Z}}
\newcommand{\Q}{\mathbb{Q}}
\newcommand{\R}{\mathbb{R}}
\newcommand{\C}{\mathbb{C}}
\newcommand{\F}{\mathbb{F}}
\newcommand{\A}{\mathbb{A}}
\newcommand{\PP}{\mathbb{P}}
\newcommand{\E}{\mathbb{E}}
\newcommand{\eps}{\varepsilon}
\newcommand{\dd}{\,\mathrm{d}}
\begin{document}
\section*{UnsolvedMath ID: 20001926}
\textbf{Source problem code:} catalogue entry
\subsection{Definitions and mathematical statement}
Let \(M\) be a closed smooth manifold and \(g\) a smooth pseudo-Riemannian metric of fixed signature \((n-q,q)\). Geodesic completeness means that every maximal affinely parameterized solution of \(\nabla^g_{\dot\gamma}\dot\gamma=0\) is defined on all of \(\R\), for every initial vector in \(TM\), including null vectors. Question 9.3.1 asks for geometric assumptions, easy to verify or satisfied by many interesting metrics, which imply this completeness. It does not ask for a specified necessary-and-sufficient classification. The recovered section introduction contrasts automatic completeness for closed Riemannian manifolds with incomplete closed examples in every indefinite signature.

The supplied research aid reports the following sufficient criterion. Suppose there is a smooth symmetric covariant 2-tensor \(K\) satisfying
\[
\nabla_iK_{jk}+\nabla_jK_{ki}+\nabla_kK_{ij}=0,
\qquad K_x(v,v)>0\quad(0\ne v\in T_xM),
\]
where \(\nabla=\nabla^g\). Then \((M,g)\) is geodesically complete. The first condition makes \(K(\dot\gamma,\dot\gamma)\) a conserved quadratic first integral; positive definiteness and compactness of \(M\) make each nonnegative level compact in \(TM\), giving the stated ODE continuation criterion. \(K\) need not be parallel or generated by Killing vector fields.

A concrete corollary in the aid assumes \(q\) global Killing fields \(X_1,\ldots,X_q\) spanning a negative-definite rank-\(q\) distribution everywhere. With \(\xi_a=g(X_a,\cdot)\), the explicit form of the coercive tensor is
\[
K_c=g+c\sum_{a=1}^q\xi_a\otimes\xi_a.
\]
For sufficiently large constant \(c>0\), it is positive definite and is a Killing tensor. For instance, writing \(A(x)=(-g(X_a,X_b))_{a,b}\), a sufficient choice is \(c>1/\min_{x\in M}\lambda_{\min}A(x)\). The minimum is positive by the stated rank, definiteness and compactness assumptions. The criterion covers nonhomogeneous torus-warped examples as reported by the aid; it does not show that all closed indefinite metrics are complete.

The original passage cites completeness of compact homogeneous pseudo-Riemannian manifolds, and of compact Lorentz manifolds of constant curvature (flat and general cases separately), and refers to S\'anchez for further results and questions. Its last sentence announces subsequent questions but does not include them. These classical cases and the tensor criterion are supplied sufficient conditions, not an exhaustive answer to the broad research question.
\subsection{Short English statement}
Find practical geometric conditions ensuring geodesic completeness of a closed metric of possibly indefinite signature. The supplied aid offers a positive-definite Killing 2-tensor as a sufficient certificate, including a Killing-field corollary that controls all velocity directions.
\subsection{Sources}
\begin{itemize}
\item Question 9.3.1, section \emph{Completeness of closed manifolds of arbitrary signature}, communicated by H. Baum, in the AIM geodesics problem list: \url{https://aimath.org/pastworkshops/geodesicsproblems.pdf}.
\item Supplied classical references: [Marsden1972] for compact homogeneous manifolds, [Carri\`ere1989] for flat compact Lorentz manifolds, [Klingler1996] for constant-curvature compact Lorentz manifolds, and [Sanchez2013] for further results and questions. Full bibliographic entries are absent.
\item Local registry \texttt{data/unsolvedmath/problems.json}, permanent ID 20001926, code \texttt{catalogue entry}; canonical \texttt{aim-geometry-notes.json}, \texttt{notes[263]}, source item 9.3; updated 14 August 2026. Clean statement null; status reconstructed unverified. The verification note recovers the section introduction from the preceding extract.
\item Research attempt 1: partial result, partially solved status, rigor label proved, novelty confidence low. The machine-generated aid is not peer reviewed and requires independent verification before citation.
\end{itemize}
\end{document}
